\section{Discussion}

The automatic score-level metrics on the logged RTX split support a bounded neural-symbolic interpretation. Harmonic conditioning and the proposed Transformer architecture improved token prediction relative to the vanilla/no-constraints Transformer, and rule-guided decoding reduced several explicit common-practice diagnostics. The most convincing improvements were observed for parallel fifths, voice crossing, spacing, and seventh-resolution violations, which are rule classes that can be identified directly from the score-level SATB representation.

The same results also show that the current method is not a complete constraint solver. Total rule violations and parallel octaves were higher for the proposed model than for the vanilla/no-constraints Transformer in the logged primary comparison. One likely reason is that local repairs can alter the texture in ways that satisfy one diagnostic while increasing another. Another is that leading-tone and cadence-related checks depend on automatically inferred harmonic context, which is necessarily uncertain in the presence of passing tones, suspensions, and modulation.

The ablations clarify the roles of the system components. Removing harmonic conditioning largely removed the likelihood advantage, which indicates that automatically derived harmonic context is useful for neural prediction. Removing iterative refinement had little effect on pitch accuracy in the logged run, but changed the distribution of perfect-interval errors. Removing rule-guided decoding kept likelihood nearly unchanged while increasing selected structural violations. These observations suggest that future work should formulate decoding as explicit multi-objective search over conflicting musical constraints rather than as a sequence of local repairs.

Human evaluation remains essential. Automatic rule reports are valuable because they make failures inspectable, but they do not measure whether a cadence sounds convincing, whether an inner voice is singable, or whether a generated solution would be useful in composition pedagogy. The prepared blind-evaluation package is therefore not a cosmetic add-on. It is needed to test whether improvements in token and rule metrics correspond to musically meaningful judgments.
